%  Copyright 2026 Sean M. Brennan and contributors
%
%   Licensed under the Apache License, Version 2.0 (the "License");
%   you may not use this file except in compliance with the License.
%   You may obtain a copy of the License at
%
%     http://www.apache.org/licenses/LICENSE-2.0
%
%   Unless required by applicable law or agreed to in writing, software
%   distributed under the License is distributed on an "AS IS" BASIS,
%   WITHOUT WARRANTIES OR CONDITIONS OF ANY KIND, either express or implied.
%   See the License for the specific language governing permissions and
%   limitations under the License.

% Passed to pandoc with --include-in-header, so this is preamble material only.

% ---- page furniture -------------------------------------------------------
\usepackage{etoolbox}
\usepackage{fancyhdr}
\usepackage{fvextra}
\usepackage{newunicodechar}
% no hyphenat/htt here on purpose. Hyphenating typewriter text inserts a
% hyphen inside an identifier, so a file name breaks as proto-cols.md and
% reads as though the hyphen were part of it. The punctuation breaks below
% do the same job without inventing a character.
\usepackage{xurl}

% Paths and dotted identifiers run through the prose and through narrow table
% cells, and TeX will not break them anywhere, so each slash and underscore is
% given a cheap break opportunity after it. Math mode is exempt.
\ifdefined\newXeTeXintercharclass
  \newXeTeXintercharclass\brkclass
  \newXeTeXintercharclass\ttbrkclass
\else
  \chardef\brkclass=200
  \chardef\ttbrkclass=201
\fi
\XeTeXcharclass`\/=\brkclass
\XeTeXcharclass`\_=\brkclass
% a dot, a colon and an open bracket separate the parts of a dotted identifier,
% but they are ordinary punctuation in prose, so these break only in typewriter
% text, where the long identifiers and the narrow table columns both live
\XeTeXcharclass`\.=\ttbrkclass
\XeTeXcharclass`\:=\ttbrkclass
\XeTeXcharclass`\(=\ttbrkclass
\makeatletter
\newcommand{\ttbrk}{\ifx\f@family\ttdefault\penalty200 \fi}
\makeatother
\XeTeXinterchartoks \brkclass 0 = {\penalty100 }
\XeTeXinterchartoks \brkclass \ttbrkclass = {\penalty100 }
\XeTeXinterchartoks \ttbrkclass 0 = {\ttbrk}
\XeTeXinterchartoks \ttbrkclass \brkclass = {\ttbrk}
\XeTeXinterchartokenstate=1
% LaTeX's own \_ carries a discretionary that breaks with a hyphen, for the
% same misreading; the break is kept and the hyphen dropped.
\DeclareTextCommand{\textunderscore}{TU}{\char`\_}
\everymath{\XeTeXinterchartokenstate=0}
\everydisplay{\XeTeXinterchartokenstate=0}

\setcounter{secnumdepth}{1}
\setlength{\emergencystretch}{4em}
\raggedbottom

% Only the per-document \markboth below may set the running head, so the
% header always names the source file rather than the nearest heading.
\renewcommand{\chaptermark}[1]{}
\renewcommand{\sectionmark}[1]{}
\setlength{\headheight}{15pt}
\pagestyle{fancy}
\fancyhf{}
\fancyhead[L]{\footnotesize\ttfamily\nouppercase{\leftmark}}
\fancyhead[R]{\footnotesize\thepage}
\renewcommand{\headrulewidth}{0.4pt}
\fancypagestyle{plain}{%
  \fancyhf{}%
  \fancyhead[L]{\footnotesize\ttfamily\nouppercase{\leftmark}}%
  \fancyhead[R]{\footnotesize\thepage}%
  \renewcommand{\headrulewidth}{0.4pt}%
}

% the source path repeated under each chapter title
\definecolor{srcgray}{gray}{0.35}
\definecolor{linkblue}{RGB}{20,60,140}
\newcommand{\srcpathline}[1]{%
  \nopagebreak\par\vspace{-0.4\baselineskip}%
  {\small\ttfamily\color{srcgray}#1}\par\vspace{0.6\baselineskip}\nopagebreak%
}

% ---- code, tables ---------------------------------------------------------
\fvset{fontsize=\small,breaklines=true,breakanywhere=true,breaksymbolleft={}}
% pandoc emits plain verbatim for unclassified code blocks, which cannot
% break a long line; route it through fancyvrb so those wrap as well
\RecustomVerbatimEnvironment{verbatim}{Verbatim}{}
\AtBeginEnvironment{longtable}{\footnotesize}
\AtBeginEnvironment{tabular}{\footnotesize}

% ---- glyph coverage -------------------------------------------------------
% TeX Gyre Pagella carries the prose; DejaVu Sans fills in the few characters
% it lacks. Verbatim material is set in DejaVu Sans Mono, which covers the
% box-drawing and geometric glyphs used in the ASCII diagrams.
\newfontfamily\fbfont{DejaVu Sans}
\newfontfamily\hebfont{DejaVu Sans}
\newcommand{\hebtext}[1]{{\hebfont #1}}
\newunicodechar{₀}{\textsubscript{0}}
\newunicodechar{₁}{\textsubscript{1}}
\newunicodechar{₂}{\textsubscript{2}}
\newunicodechar{ἀ}{{\fbfont ἀ}}
\newunicodechar{ἔ}{{\fbfont ἔ}}
\newunicodechar{ά}{{\fbfont ά}}
\newunicodechar{ί}{{\fbfont ί}}
\newunicodechar{►}{{\fbfont ►}}
\newunicodechar{☆}{{\fbfont ☆}}
\newunicodechar{⟂}{\ensuremath{\perp}}
